\documentclass[11pt]{article}
\usepackage[margin=0.9in]{geometry}
\usepackage{amsmath,amssymb,booktabs,enumitem,listings,xcolor,hyperref}
\definecolor{backcolour}{rgb}{0.95,0.95,0.92}
\definecolor{codeblue}{rgb}{0.05,0.15,0.55}
\definecolor{codegreen}{rgb}{0.0,0.45,0.0}
\lstset{
    backgroundcolor=\color{backcolour},
    basicstyle=\ttfamily\small,
    keywordstyle=\color{codeblue}\bfseries,
    commentstyle=\color{codegreen},
    stringstyle=\color{purple},
    breaklines=true,
    frame=single,
    showstringspaces=false,
    tabsize=4
}

\title{CPSC 250 \\ Class Constructors}
\author{Edward Brash}
\date{}
\begin{document}
\maketitle

\section{Review}

In the previous class examples, the constructor was kept as simple as possible.

For example:

\begin{lstlisting}[language=Python]
def __init__(self):
    self.name = "none"
    self.rating = -1
    self.price = "none"
    self.cuisine = "none"
\end{lstlisting}

This constructor creates an object and assigns default values to all member variables.

That works, but it is not always convenient.

\section{The Problem}

Often, when we create an object, we already know some or all of the information that should go inside it.

For example, rather than writing:

\begin{lstlisting}[language=Python]
moes = Restaurant()

moes.set_name("Moe's")
moes.set_rating(3)
moes.set_price("***")
moes.set_cuisine("Mexican")
\end{lstlisting}

it would be nicer to write:

\begin{lstlisting}[language=Python]
moes = Restaurant("Moe's", 3, "***", "Mexican")
\end{lstlisting}

This is possible by adding parameters to the constructor.

\section{Constructor Parameters}

We can define the constructor as:

\begin{lstlisting}[language=Python]
class Restaurant:

    def __init__(self, name, rating, price, cuisine):
        self.name = name
        self.rating = rating
        self.price = price
        self.cuisine = cuisine
\end{lstlisting}

Now the object can be initialized directly:

\begin{lstlisting}[language=Python]
moes = Restaurant("Moe's", 3, "***", "Mexican")
\end{lstlisting}

\section{Important Notes}

\subsection{Parameter Names}

It is common to use the same names for constructor parameters and instance attributes:

\begin{lstlisting}[language=Python]
def __init__(self, name, rating, price, cuisine):
    self.name = name
    self.rating = rating
    self.price = price
    self.cuisine = cuisine
\end{lstlisting}

This is not a conflict.

The name on the right side is the parameter. The name on the left side after \verb|self.| is the instance attribute.

\subsection{Documentation}

It is good practice to include documentation explaining the parameters.

\begin{lstlisting}[language=Python]
class Restaurant:
    # Represents a restaurant.
    #
    # Parameters:
    #     name: Name of the restaurant.
    #     rating: Rating from 1 to 10.
    #     price: Price category such as "$", "$$", or "$$$".
    #     cuisine: Type of cuisine.
\end{lstlisting}

This documentation can be viewed using:

\begin{lstlisting}[language=Python]
help(Restaurant)
\end{lstlisting}

\section{Default Constructor Values}

Sometimes we want to provide default values for some constructor arguments.

Example:

\begin{lstlisting}[language=Python]
class Restaurant:

    def __init__(self, name, cuisine, rating=-1, price="none"):
        self.name = name
        self.cuisine = cuisine
        self.rating = rating
        self.price = price
\end{lstlisting}

Now we can write:

\begin{lstlisting}[language=Python]
moes = Restaurant("Moe's", "Mexican")
\end{lstlisting}

and the object will use default values for \verb|rating| and \verb|price|.

\section{Important Rule About Defaults}

Non-default parameters must come before default parameters.

This is valid:

\begin{lstlisting}[language=Python]
def __init__(self, name, cuisine, rating=-1, price="none"):
\end{lstlisting}

This is not valid:

\begin{lstlisting}[language=Python]
def __init__(self, name="none", cuisine, rating=-1):
\end{lstlisting}

Python needs all required parameters first, followed by optional parameters.

\section{Overriding Default Values}

Default values are convenient, but they can still be overridden.

\begin{lstlisting}[language=Python]
panera = Restaurant("Panera", "Soups and Sandwiches", 4, "$$$")
five_guys = Restaurant("Five Guys", "Burgers", 10, "$$$$$")
\end{lstlisting}

These calls supply explicit values for all fields.

\section{Best Practice: Give Defaults to All Member Variables}

A common design choice is to make sure every member variable has a reasonable value after construction.

This prevents partially initialized objects.

For example:

\begin{lstlisting}[language=Python]
class Restaurant:

    def __init__(self, name="none", cuisine="none",
                 rating=-1, price="none"):
        self.name = name
        self.cuisine = cuisine
        self.rating = rating
        self.price = price
\end{lstlisting}

Then all of the following make sense:

\begin{lstlisting}[language=Python]
Restaurant()
Restaurant("Moe's")
Restaurant("Moe's", "Fake Mex")
Restaurant("Moe's", "Fake Mex", 6)
Restaurant("Moe's", "Fake Mex", 6, "$$")
\end{lstlisting}

\section{Internal Helper Methods}

Another best practice is to use an underscore to indicate methods that are intended to be internal.

For example, suppose we define a rectangle:

\begin{lstlisting}[language=Python]
class Rectangle:

    def __init__(self, length=1, width=1):
        self.length = length
        self.width = width

    def _get_area(self):
        return self.length * self.width

    def resize(self, new_length):
        area = self._get_area()
        self.length = new_length
        self.width = area / new_length

    def print_rectangle(self):
        print(f"Length: {self.length:.3f}")
        print(f"Width: {self.width:.3f}")
        print(f"Area: {self._get_area():.3f}")
\end{lstlisting}

The method \verb|_get_area()| is still technically accessible from outside the class, but the underscore indicates:

\begin{quote}
This method is intended for internal use.
\end{quote}

\section{A Note About Naming Conventions}

When we see a method beginning with an underscore, such as:

\begin{lstlisting}[language=Python]
_get_area()
\end{lstlisting}

we should not assume it is magical or special.

It is mostly a convention.

Python programmers use this convention to communicate design intent.

\section{Key Takeaways}

\begin{itemize}
    \item Constructors initialize objects.
    \item Constructors can take parameters.
    \item Parameter values can be assigned to instance attributes.
    \item Default parameter values make object creation more flexible.
    \item Required parameters must come before default parameters.
    \item Documentation helps users understand constructor parameters.
    \item Leading underscores indicate internal-use methods by convention.
\end{itemize}

\end{document}
